\documentclass[11pt,a4paper]{article}

\usepackage[utf8]{inputenc}
\usepackage[T1]{fontenc}
\usepackage{amsmath,amssymb,amsthm}
\usepackage{mathtools}
\usepackage[margin=2.5cm]{geometry}
\usepackage{hyperref}

% -------------------- Environments --------------------
\newtheorem{proposition}{Proposition}[section]

% -------------------- Macros --------------------
\newcommand{\dd}{\mathrm{d}}
\newcommand{\exc}{e_{xc}}
\newcommand{\Exc}{E_{xc}}
\newcommand{\vrho}[1]{v_{\rho#1}}
\newcommand{\vsig}[1]{v_{\sigma#1}}

\title{Derivation of the weak-form GGA exchange--correlation potential\\
matrix used by \texttt{assemble\_exc!}}
\author{}
\date{\today}

\begin{document}
\maketitle

\begin{abstract}
This note derives, from the radial finite-element density parametrization
used throughout \texttt{AtomicKohnSham}, the closed-form weak-form
expression for the GGA exchange--correlation potential matrix
$V_{xc\sigma} = \partial \Exc/\partial D^\sigma$ implemented in
\texttt{assemble\_exc!} (see \texttt{dev/gga/PLAN.md} for the original,
more compressed derivation). Every factor of $2$ and $4$ is re-derived
explicitly, since this is exactly the class of bookkeeping error that is
easy to get wrong and hard to notice (a discrepancy here degrades accuracy
smoothly rather than crashing).
\end{abstract}

\section{Setup and notation}

For a fixed angular-momentum channel and spin $\sigma \in \{\uparrow,\downarrow\}$,
the radial density matrix $D^\sigma$ (symmetric, $D^\sigma_{ij} = D^\sigma_{ji}$)
parametrizes the reduced radial density through the finite-element basis
$(\varphi_i)_i$ of the functions $u(r) = rR(r)$:
\begin{equation}
    q^\sigma(r) = \sum_{ij} D^\sigma_{ij}\,\varphi_i(r)\varphi_j(r), \qquad
    \rho^\sigma(r) = \frac{q^\sigma(r)}{4\pi r^2}.
\end{equation}
The spin-polarized GGA exchange--correlation energy is
\begin{equation}
    \Exc = \int_0^\infty \exc\big(\rho_\uparrow,\rho_\downarrow,
        \sigma_{\uparrow\uparrow},\sigma_{\uparrow\downarrow},\sigma_{\downarrow\downarrow}\big)\,
        4\pi r^2\,\dd r, \qquad
    \sigma_{\alpha\beta}(r) := \rho_\alpha'(r)\,\rho_\beta'(r),
\end{equation}
where $\exc$ is the (Libxc-convention) exchange--correlation energy density,
$\vrho{\alpha} := \partial \exc/\partial\rho_\alpha$ and
$\vsig{\alpha\beta} := \partial \exc/\partial\sigma_{\alpha\beta}$.

The Kohn--Sham weak form requires the Hamiltonian matrix element
contributed by $\Exc$, obtained by differentiating $\Exc$ with respect to
each entry of $D^\uparrow$ (the $\downarrow$ channel follows by symmetry):
\begin{equation}
    V_{xc\uparrow} := \frac{\partial \Exc}{\partial D^\uparrow}.
\end{equation}
Since the collinear approximation keeps the two spin channels' density
matrices independent, $\rho_\downarrow$ and $\sigma_{\downarrow\downarrow}$
do not depend on $D^\uparrow$, and only three chain-rule terms survive.

\section{Step 1: the density gradient as a linear map of \texorpdfstring{$D$}{D}}

By the quotient rule, $\rho' = q'/(4\pi r^2) - 2\rho/r$. Because $D$ is
symmetric, relabeling the dummy indices $i \leftrightarrow j$ in half of
the product-rule expansion collapses the sum:
\begin{equation}
    q'(r) = \sum_{ij} D_{ij}\big(\varphi_i'\varphi_j + \varphi_i\varphi_j'\big)
          = 2\sum_{ij} D_{ij}\,\varphi_i'\varphi_j .
    \label{eq:qprime}
\end{equation}
Differentiating $\rho'$ with respect to a single (unsymmetrized) entry
$D_{ij}$ -- the mass-matrix building blocks introduced in Step~3 restore
the correct symmetric matrix, so no care is needed here about
double-counting off-diagonal entries -- gives, using \eqref{eq:qprime} and
the $-2\rho/r$ term (which itself has $\partial\rho/\partial D_{ij} =
\varphi_i\varphi_j/4\pi r^2$):
\begin{equation}
    \frac{\partial \rho'(r)}{\partial D_{ij}}
    = \frac{1}{4\pi r^2}\left[\varphi_i'\varphi_j + \varphi_i\varphi_j'
        - \frac{2}{r}\varphi_i\varphi_j\right].
    \label{eq:drhoprime}
\end{equation}
The $-2/r$ term is a purely radial-reduction artifact: it comes from
$q = 4\pi r^2\rho$, not from any physical content of the GGA functional,
and would be absent in a Cartesian weak-form derivation.

\section{Step 2: chain rule through \texorpdfstring{$\rho_\uparrow$, $\sigma_{\uparrow\uparrow}$, $\sigma_{\uparrow\downarrow}$}{rho, sigma-up-up, sigma-up-down}}

\begin{equation}
    \frac{\partial \Exc}{\partial D^\uparrow_{ij}}
    = \int_0^\infty 4\pi r^2\left[
        \vrho{\uparrow}\,\frac{\partial \rho_\uparrow}{\partial D^\uparrow_{ij}}
        + \vsig{\uparrow\uparrow}\,\frac{\partial \sigma_{\uparrow\uparrow}}{\partial D^\uparrow_{ij}}
        + \vsig{\uparrow\downarrow}\,\frac{\partial \sigma_{\uparrow\downarrow}}{\partial D^\uparrow_{ij}}
      \right]\dd r.
\end{equation}
Using $\sigma_{\uparrow\uparrow} = \rho_\uparrow'^2$ and
$\sigma_{\uparrow\downarrow} = \rho_\uparrow'\rho_\downarrow'$ (product rule):
\begin{equation}
    \frac{\partial \sigma_{\uparrow\uparrow}}{\partial D^\uparrow_{ij}}
        = 2\rho_\uparrow'\,\frac{\partial \rho_\uparrow'}{\partial D^\uparrow_{ij}},
    \qquad
    \frac{\partial \sigma_{\uparrow\downarrow}}{\partial D^\uparrow_{ij}}
        = \rho_\downarrow'\,\frac{\partial \rho_\uparrow'}{\partial D^\uparrow_{ij}}.
\end{equation}
Both multiply the \emph{same} factor $\partial\rho_\uparrow'/\partial D^\uparrow_{ij}$
from \eqref{eq:drhoprime}, so they combine into a single coefficient:
\begin{equation}
    K_\uparrow(r) := 2\vsig{\uparrow\uparrow}(r)\,\rho_\uparrow'(r)
        + \vsig{\uparrow\downarrow}(r)\,\rho_\downarrow'(r).
    \label{eq:Kup}
\end{equation}
Substituting \eqref{eq:drhoprime}--\eqref{eq:Kup} and
$\partial\rho_\uparrow/\partial D^\uparrow_{ij} = \varphi_i\varphi_j/4\pi r^2$,
the factors of $4\pi r^2$ cancel exactly:
\begin{align}
    \frac{\partial \Exc}{\partial D^\uparrow_{ij}}
    &= \int_0^\infty \left\{ \vrho{\uparrow}\,\varphi_i\varphi_j
        + K_\uparrow\left[\varphi_i'\varphi_j + \varphi_i\varphi_j' - \frac{2}{r}\varphi_i\varphi_j\right]\right\}\dd r \\
    &= \int_0^\infty \Big[\vrho{\uparrow} - \tfrac{2}{r}K_\uparrow\Big]\varphi_i\varphi_j\,\dd r
        + \int_0^\infty K_\uparrow\,\varphi_i'\varphi_j\,\dd r
        + \int_0^\infty K_\uparrow\,\varphi_i\varphi_j'\,\dd r.
    \label{eq:threeterms}
\end{align}

\section{Step 3: identifying the matrix building blocks}

Define the weighted mass matrix and the (asymmetric) mixed mass matrix
already used elsewhere in the codebase for the LDA and Hartree assembly:
\begin{equation}
    M[w]_{ij} := \int_0^\infty w(r)\,\varphi_i(r)\varphi_j(r)\,\dd r
    \quad(\text{symmetric}),\qquad
    Mmix[w]_{ij} := \int_0^\infty w(r)\,\varphi_i'(r)\varphi_j(r)\,\dd r
    \quad(\text{not symmetric}).
\end{equation}
$M[w]$ is symmetric because $\varphi_i\varphi_j$ does not distinguish the
order of $i,j$; $Mmix[w]$ is not, since $\varphi_i'\varphi_j \ne \varphi_j'\varphi_i$
in general. The third integral in \eqref{eq:threeterms} is exactly the
transpose of the second under this definition:
\begin{equation}
    \int_0^\infty K_\uparrow\,\varphi_i\varphi_j'\,\dd r
    = \int_0^\infty K_\uparrow\,\varphi_j'\varphi_i\,\dd r
    = Mmix[K_\uparrow]_{ji} = \big(Mmix[K_\uparrow]^T\big)_{ij}.
\end{equation}
Hence, in matrix form:
\begin{equation}
    \boxed{V_{xc\uparrow} = M\Big[\vrho{\uparrow} - \frac{2}{r}K_\uparrow\Big]
        + Mmix[K_\uparrow] + Mmix[K_\uparrow]^T}, \qquad
    K_\uparrow = 2\vsig{\uparrow\uparrow}\rho_\uparrow' + \vsig{\uparrow\downarrow}\rho_\downarrow'
    \label{eq:final}
\end{equation}
and $V_{xc\downarrow}$, $K_\downarrow$ follow by exchanging
$\uparrow\leftrightarrow\downarrow$ throughout. The
$Mmix[K_\uparrow] + Mmix[K_\uparrow]^T$ combination is what makes
$V_{xc\uparrow}$ symmetric despite each piece individually not being so --
physically required of any Hamiltonian matrix, and structurally the same
defensive $(V+V^T)/2$ pattern already used in \texttt{assemble\_exc!}.

\begin{proposition}[Consistency check: $\nabla$-independent limit]
Setting $\vsig{\alpha\beta} \equiv 0$ (an LDA functional has no gradient
dependence) collapses $K_\uparrow \to 0$ and \eqref{eq:final} reduces to
$V_{xc\uparrow} = M[\vrho{\uparrow}]$, the ordinary LDA weighted mass
matrix -- as it must.
\end{proposition}

\section{Step 4: the unpolarized case, \texorpdfstring{$n_{spin}=1$}{nspin=1}}

An unpolarized GGA functional depends on a single gradient invariant
$\sigma = \rho'^2$ rather than the three spin-resolved invariants above --
there is no second channel to mix with, so no cross term arises. Repeating
Steps~1--3 with $\sigma_{\uparrow\uparrow} \to \sigma$,
$\sigma_{\uparrow\downarrow} \to 0$ gives simply
\begin{equation}
    K = 2v_\sigma\,\rho', \qquad
    V_{xc} = M\Big[\vrho{} - \frac{2}{r}K\Big] + Mmix[K] + Mmix[K]^T
           = M\Big[\vrho{} - \frac{4}{r}v_\sigma\rho'\Big]
             + Mmix[2v_\sigma\rho'] + Mmix[2v_\sigma\rho']^T,
\end{equation}
matching the single-channel formula originally derived in
\texttt{dev/gga/PLAN.md}.

\end{document}
